% ==============================================================================
% ภาคผนวก (Appendix)
% โครงการ: การพัฒนาแพลตฟอร์มสำหรับการสนับสนุนการวินิจฉัยรอยโรคทางช่องปาก
%          และผสมผสานการใช้งานระบบผู้เชี่ยวชาญสำหรับทันตแพทย์
% ==============================================================================

% ------------------------------------------------------------------------------
% ภาคผนวก ก: คู่มือการติดตั้งและกำหนดค่าสภาพแวดล้อมระบบ
% ------------------------------------------------------------------------------
\chapter{คู่มือการติดตั้งและการกำหนดค่าสภาพแวดล้อมระบบ}
\label{appendix:installation}

ภาคผนวกนี้อธิบายข้อกำหนดของระบบ โครงสร้างไฟล์ และขั้นตอนการติดตั้งสภาพแวดล้อมการทำงานของระบบสนับสนุนการวินิจฉัยรอยโรคทางช่องปาก โดยใช้เทคโนโลยีคอนเทนเนอร์ (Docker และ Docker Compose) เพื่ออำนวยความสะดวกในการนำระบบไปติดตั้งและใช้งานจริง

\section{ข้อกำหนดของระบบ (System Requirements)}
เพื่อให้ระบบสามารถประมวลผลแบบจำลองปัญญาประดิษฐ์และให้บริการเว็บแอปพลิเคชันได้อย่างมีประสิทธิภาพ อุปกรณ์เซิร์ฟเวอร์ควรมีคุณสมบัติขั้นต่ำดังนี้:

\begin{enumerate}
    \item \textbf{ข้อกำหนดด้านฮาร์ดแวร์ (Hardware Requirements):}
    \begin{itemize}
        \item หน่วยประมวลผลกลาง (CPU): สถาปัตยกรรม x86\_64 ความเร็วไม่น้อยกว่า 2.0 GHz จำนวน 4 แกนประมวลผล (Cores) ขึ้นไป
        \item หน่วยความจำหลัก (RAM): ไม่น้อยกว่า 8 GB (แนะนำ 16 GB ขึ้นไป สำหรับการรันแบบจำลอง Multimodal คู่ขนาน)
        \item พื้นที่จัดเก็บข้อมูล (Storage): ไม่น้อยกว่า 15 GB (สำหรับจัดเก็บ Docker Images, Weights ของโมเดล และไฟล์ภาพประวัติ)
        \item หน่วยประมวลผลกราฟิก (GPU - ทางเลือก): NVIDIA GPU รองรับ CUDA 11.8+ พร้อม VRAM ไม่น้อยกว่า 6 GB (หากไม่มี ระบบจะสลับไปประมวลผลบน CPU อัตโนมัติ)
    \end{itemize}
    \item \textbf{ข้อกำหนดด้านซอฟต์แวร์ (Software Requirements):}
    \begin{itemize}
        \item ระบบปฏิบัติการ: Ubuntu Linux 22.04 LTS, Debian 12 หรือ Windows 11 ที่เปิดใช้งาน WSL2 (Windows Subsystem for Linux)
        \item ซอฟต์แวร์คอนเทนเนอร์: Docker Engine เวอร์ชัน 24.0.0 ขึ้นไป และ Docker Compose Plugin เวอร์ชัน 2.20.0 ขึ้นไป
        \item เว็บบราวเซอร์สำหรับผู้ใช้งาน: Google Chrome, Mozilla Firefox หรือ Microsoft Edge เวอร์ชันล่าสุด
    \end{itemize}
\end{enumerate}

\section{โครงสร้างไดเรกทอรีของโครงงาน (Project Directory Structure)}
โครงสร้างแฟ้มข้อมูลหลักของระบบสนับสนุนการวินิจฉัยรอยโรคทางช่องปาก ถูกจัดระเบียบตามสถาปัตยกรรมไมโครเซอร์วิส ดังนี้:

\begin{verbatim}
Oral_Lesion_2026/
|-- docker-compose.yml              # ไฟล์คอนฟิกูเรชันสำหรับควบคุมทุกคอนเทนเนอร์
|-- test_api.py                     # สคริปต์ทดสอบการทำงานของ API อัตโนมัติ
|-- frontend/                       # ไมโครเซอร์วิสส่วนติดต่อผู้ใช้งาน (React + Vite)
|   |-- Dockerfile
|   |-- package.json
|   |-- src/
|   `-- tailwind.config.js
|-- backend/                        # ไมโครเซอร์วิส Backend API Gateway (FastAPI)
|   |-- Dockerfile
|   |-- requirements.txt
|   |-- main.py
|   `-- static/
|       |-- uploads/                # ไดเรกทอรีเก็บภาพที่อัปโหลดเข้าสู่ระบบ
|       `-- outputs/                # ไดเรกทอรีเก็บภาพผลลัพธ์ (BBox / Grad-CAM)
`-- ai-engine/                      # ไมโครเซอร์วิสประมวลผลปัญญาประดิษฐ์
    |-- Dockerfile
    |-- requirements.txt
    |-- main.py
    |-- yolo-service/               # โมดูลตรวจจับวัตถุ YOLOv7 และ YOLOv8
    |   |-- yolov7.pt
    |   `-- yolov8.pt
    `-- densenet-bert-service/      # โมดูลการเรียนรู้เชิงลึกแบบหลายโมดาลิตี
        |-- best_densenet121_img.pth
        `-- best_text.pth
\end{verbatim}

\section{การจัดเตรียมไฟล์น้ำหนักแบบจำลอง (Model Weights Setup)}
ก่อนเริ่มต้นการทำงานของระบบ ต้องตรวจสอบให้แน่ใจว่าได้วางไฟล์น้ำหนักแบบจำลอง (Pre-trained Weights) ไว้ในตำแหน่งที่ถูกต้อง ดังนี้:
\begin{enumerate}
    \item ไฟล์แบบจำลอง YOLOv7: จัดเก็บที่ \texttt{ai-engine/yolo-service/yolov7.pt}
    \item ไฟล์แบบจำลอง YOLOv8: จัดเก็บที่ \texttt{ai-engine/yolo-service/yolov8.pt}
    \item ไฟล์แบบจำลอง DenseNet-121 สำหรับประมวลผลภาพ: จัดเก็บที่ \texttt{ai-engine/densenet-bert-service/best\_densenet121\_img.pth}
    \item ไฟล์แบบจำลอง WangchanBERTa สำหรับประมวลผลข้อความ: จัดเก็บที่ \texttt{ai-engine/densenet-bert-service/best\_text.pth}
\end{enumerate}

\section{การกำหนดค่าตัวแปรสภาพแวดล้อม (Environment Variables)}
ระบบควบคุมการทำงานของแต่ละไมโครเซอร์วิสผ่านตัวแปรสภาพแวดล้อมในไฟล์ \texttt{docker-compose.yml} ดังแสดงในตารางที่ \ref{tab:env_vars}:

\begin{table}[htbp]
\centering
\small
\caption{รายละเอียดตัวแปรสภาพแวดล้อมของระบบ}
\label{tab:env_vars}
\begin{tabular}{p{2.5cm}p{3.5cm}p{3.5cm}p{4cm}}
\toprule
\textbf{เซอร์วิส} & \textbf{ชื่อตัวแปร} & \textbf{ค่าเริ่มต้น} & \textbf{คำอธิบายหน้าที่} \\
\midrule
Frontend & \texttt{VITE\_API\_URL} & \texttt{http://localhost:8000} & ที่อยู่ URL สำหรับเชื่อมต่อ Backend Gateway \\
\midrule
Backend & \texttt{AI\_ENGINE\_URL} & \texttt{http://ai-engine:8001} & ที่อยู่ภายในวงเครือข่ายสำหรับเรียก AI Engine \\
\midrule
\multirow{2}{*}{AI Engine} & \texttt{PYTHONUNBUFFERED} & \texttt{1} & บังคับให้ Python แสดงผล log ออกทันที \\
 & \texttt{FUSION\_ALPHA} & \texttt{0.6} & ค่าน้ำหนักถ่วงภาพในการคำนวณ Late Fusion \\
\bottomrule
\end{tabular}
\end{table}

\section{ขั้นตอนการสร้างและเริ่มต้นระบบ (Deployment Steps)}
ผู้ดูแลระบบสามารถสั่งสร้างและเปิดให้บริการระบบทั้งหมดผ่านคำสั่ง Docker Compose ตามขั้นตอนดังนี้:

\begin{enumerate}
    \item เปิดเทอร์มินัลและเข้าไปยังไดเรกทอรีรากของโครงงาน:
\begin{verbatim}
cd Oral_Lesion_2026
\end{verbatim}

    \item สั่งสร้างคอนเทนเนอร์อิมเมจและเริ่มต้นระบบในโหมดเบื้องหลัง (Detached Mode):
\begin{verbatim}
docker compose up -d --build
\end{verbatim}

    \item ตรวจสอบสถานะความพร้อมของคอนเทนเนอร์ทั้ง 3 เซอร์วิส:
\begin{verbatim}
docker compose ps
\end{verbatim}
    เมื่อระบบพร้อมทำงาน คอนเทนเนอร์ \texttt{oralscan-frontend}, \texttt{oralscan-backend}, และ \texttt{oralscan-ai-engine} จะต้องแสดงสถานะเป็น \texttt{Up (healthy)} หรือ \texttt{Up}

    \item การเข้าใช้งานระบบผ่านเว็บบราวเซอร์:
    \begin{itemize}
        \item ส่วนติดต่อผู้ใช้งานสำหรับทันตแพทย์: \url{http://localhost:5173}
        \item เอกสาร API ทางการ (Swagger UI): \url{http://localhost:8000/docs}
    \end{itemize}

    \item คำสั่งสำหรับหยุดและปิดการทำงานของระบบ:
\begin{verbatim}
docker compose down
\end{verbatim}
\end{enumerate}

% ------------------------------------------------------------------------------
% ภาคผนวก ข: ข้อกำหนดส่วนต่อประสานโปรแกรมประยุกต์
% ------------------------------------------------------------------------------
\chapter{ข้อกำหนดส่วนต่อประสานโปรแกรมประยุกต์ (API Specification)}
\label{appendix:api_spec}

ภาคผนวกนี้อธิบายข้อกำหนดทางเทคนิคของ RESTful API ที่ให้บริการโดย Backend API Gateway สำหรับการเชื่อมต่อระหว่างระบบส่วนหน้า (Frontend) หรือระบบภายนอกกับระบบประเมินรอยโรคช่องปาก

\section{ภาพรวมของส่วนต่อประสาน (API Overview)}
ระบบให้บริการ API ภายใต้โปรโตคอล HTTP/1.1 และส่งคืนข้อมูลในรูปแบบ JSON (JavaScript Object Notation) โดยมี Endpoint หลักแสดงในตารางที่ \ref{tab:api_endpoints}:

\begin{table}[htbp]
\centering
\small
\caption{รายการ API Endpoints ของระบบสนับสนุนการวินิจฉัยรอยโรคทางช่องปาก}
\label{tab:api_endpoints}
\begin{tabular}{p{1.8cm}p{4.2cm}p{3.5cm}p{4.2cm}}
\toprule
\textbf{Method} & \textbf{Endpoint Path} & \textbf{Content-Type} & \textbf{คำอธิบายหน้าที่} \\
\midrule
\texttt{GET} & \texttt{/api/health} & \texttt{application/json} & ตรวจสอบสถานะความพร้อมของเซิร์ฟเวอร์ \\
\texttt{POST} & \texttt{/api/assess} & \texttt{multipart/form-data} & ส่งภาพถ่ายและข้อมูลอาการเพื่อประเมินรอยโรค \\
\texttt{GET} & \texttt{/static/outputs/\{file\}} & \texttt{image/jpeg} หรือ \texttt{image/png} & ให้บริการเรียกดูไฟล์ภาพผลลัพธ์ \\
\bottomrule
\end{tabular}
\end{table}

\section{ข้อกำหนดของ Endpoint ประเมินรอยโรค (/api/assess)}

\subsection{ข้อมูลนำเข้า (Request Parameters)}
คำขอส่งตรวจประเมินรอยโรคต้องส่งด้วยวิธี \texttt{POST} ในรูปแบบ \texttt{multipart/form-data} ประกอบด้วยฟิลด์ข้อมูลดังนี้:

\begin{enumerate}
    \item \textbf{\texttt{image} (จำเป็น):} ไฟล์ภาพถ่ายรอยโรคในช่องปาก รองรับนามสกุลไฟล์ \texttt{.jpg}, \texttt{.jpeg}, และ \texttt{.png} ขนาดไฟล์ไม่เกิน 10 MB
    \item \textbf{\texttt{modelType} (จำเป็น):} ข้อความระบุชื่อสถาปัตยกรรมแบบจำลองที่ต้องการใช้ประมวลผล รองรับ 3 ค่า ได้แก่:
    \begin{itemize}
        \item \texttt{"YOLOv7"} สำหรับตรวจจับตำแหน่งด้วยแบบจำลอง YOLOv7
        \item \texttt{"YOLOv8"} สำหรับตรวจจับตำแหน่งด้วยแบบจำลอง YOLOv8
        \item \texttt{"DenseNet-BERT"} สำหรับการจำแนกหลายโมดาลิตีและสร้างแผนที่ Grad-CAM Heatmap
    \end{itemize}
    \item \textbf{\texttt{text} (ทางเลือก):} ข้อความบรรยายอาการ พฤติกรรมเสี่ยง หรือคำสำคัญทางคลินิก (เช่น \texttt{"smoking, ulcer, ดื่มเหล้า"}) ซึ่งจำเป็นต้องระบุหากเลือกใช้งานโมเดล \texttt{"DenseNet-BERT"} เพื่อให้โมดูล WangchanBERTa นำไปประมวลผล
\end{enumerate}

\subsection{โครงสร้างข้อมูลตอบกลับสำเร็จ (Success Response: HTTP 200 OK)}
เมื่อระบบประมวลผลคำขอสำเร็จ จะส่งคืนข้อมูล JSON ตามโครงสร้างตัวอย่างดังนี้:

\begin{verbatim}
{
  "success": true,
  "model_used": "YOLOv7",
  "result": {
    "disease_name": "Leukoplakia",
    "risk_level": "High",
    "confidence_score": 0.884,
    "bounding_boxes": [
      {
        "xmin": 142.5,
        "ymin": 210.0,
        "xmax": 358.2,
        "ymax": 420.7,
        "label": "Leukoplakia",
        "conf": 0.884
      }
    ],
    "processed_image_url": "/static/outputs/annotated_uuid123.jpg"
  }
}
\end{verbatim}

\subsection{รหัสสถานะและการจัดการข้อผิดพลาด (Error Handling)}
ระบบมีการตรวจสอบความถูกต้องของข้อมูลนำเข้าและส่งคืนรหัสสถานะ HTTP ตามมาตรฐานสากล:
\begin{itemize}
    \item \textbf{\texttt{400 Bad Request}:} เกิดขึ้นเมื่อข้อมูลนำเข้าไม่ถูกต้อง เช่น ไม่อัปโหลดไฟล์ภาพ, นามสกุลไฟล์ไม่ใช่ภาพ, หรือระบุชื่อ \texttt{modelType} ไม่ตรงกับที่ระบบรองรับ
    \item \textbf{\texttt{500 Internal Server Error}:} เกิดขึ้นเมื่อโมดูล AI Engine ขัดข้อง หรือหน่วยความจำไม่เพียงพอขณะทำการอนุมานผล
\end{itemize}

% ------------------------------------------------------------------------------
% ภาคผนวก ค: สคริปต์การทดสอบและผลการทดสอบระบบอัตโนมัติ
% ------------------------------------------------------------------------------
\chapter{สคริปต์การทดสอบและผลการทดสอบระบบอัตโนมัติ}
\label{appendix:test_scripts}

ภาคผนวกนี้แสดงซอร์สโค้ดของสคริปต์การทดสอบระบบอัตโนมัติ (\texttt{test\_api.py}) ซึ่งใช้สำหรับตรวจสอบการทำงานของ Backend API และตัวอย่างผลลัพธ์การตอบกลับที่ได้รับจากการทดสอบจริง

\section{ซอร์สโค้ดสคริปต์การทดสอบ (test\_api.py)}
สคริปต์ทดสอบถูกพัฒนาขึ้นด้วยภาษา Python โดยใช้ไลบรารีมาตรฐาน \texttt{urllib.request} เพื่อจำลองการส่งคำขอแบบ \texttt{multipart/form-data} ไปยัง Endpoint \texttt{/api/assess}:

\begin{verbatim}
import urllib.request
import json
from pathlib import Path

def test_endpoint(model_type, text=""):
    url = "http://localhost:8000/api/assess"
    boundary = "----WebKitFormBoundaryTest123"
    
    parts = []
    # 1. ฟิลด์ระบุประเภทโมเดล (modelType)
    parts.append(
        f"--{boundary}\r\n"
        f"Content-Disposition: form-data; name=\"modelType\"\r\n\r\n"
        f"{model_type}\r\n".encode("utf-8")
    )
    
    # 2. ฟิลด์ระบุข้อความอาการ (text)
    if text:
        parts.append(
            f"--{boundary}\r\n"
            f"Content-Disposition: form-data; name=\"text\"\r\n\r\n"
            f"{text}\r\n".encode("utf-8")
        )
        
    # 3. แนบไฟล์ภาพตัวอย่าง (image)
    parts.append(
        f"--{boundary}\r\n"
        f"Content-Disposition: form-data; name=\"image\"; filename=\"test.jpg\"\r\n"
        f"Content-Type: image/jpeg\r\n\r\n".encode("utf-8")
    )
    with open("static/uploads/test.jpg", "rb") as f:
        parts.append(f.read())
    parts.append(f"\r\n--{boundary}--\r\n".encode("utf-8"))
    
    body = b"".join(parts)
    req = urllib.request.Request(
        url,
        data=body,
        headers={"Content-Type": f"multipart/form-data; boundary={boundary}"}
    )
    
    with urllib.request.urlopen(req) as resp:
        res = json.loads(resp.read().decode("utf-8"))
        print(f"\n[+] Test {model_type} -> HTTP {resp.status}")
        print(json.dumps(res, indent=2, ensure_ascii=False))

if __name__ == "__main__":
    print("Testing YOLOv7 Model...")
    test_endpoint("YOLOv7", "smoking")
    
    print("\nTesting YOLOv8 Model...")
    test_endpoint("YOLOv8", "alwaysHurts")
    
    print("\nTesting DenseNet-BERT Model...")
    test_endpoint("DenseNet-BERT", "ดื่มเครื่องดื่มแอลกอฮอล์ สูบบุหรี่ ปื้นขาว ลิ้น")
\end{verbatim}

\section{ตัวอย่างผลลัพธ์จากการทดสอบแบบจำลองแต่ละประเภท}

\subsection{ผลการทดสอบแบบจำลอง YOLOv7}
ผลลัพธ์การตอบกลับเมื่อส่งทดสอบภาพถ่ายรอยโรคด้วยโมเดล YOLOv7 พร้อมคำสำคัญ \texttt{"smoking"}:

\begin{verbatim}
[+] Test YOLOv7 -> HTTP 200
{
  "success": true,
  "model_used": "YOLOv7",
  "result": {
    "disease_name": "Leukoplakia",
    "risk_level": "High",
    "confidence_score": 0.884,
    "bounding_boxes": [
      {
        "xmin": 154.2,
        "ymin": 204.6,
        "xmax": 372.8,
        "ymax": 418.1,
        "label": "Leukoplakia",
        "conf": 0.884
      }
    ],
    "processed_image_url": "/static/outputs/annotated_yolo7_f8a2c1.jpg"
  }
}
\end{verbatim}

\subsection{ผลการทดสอบแบบจำลอง Multimodal DenseNet-BERT}
ผลลัพธ์การตอบกลับเมื่อส่งทดสอบภาพถ่ายรอยโรคด้วยโมเดล Multimodal DenseNet-BERT ร่วมกับข้อความอาการภาษาไทย \texttt{"ดื่มเครื่องดื่มแอลกอฮอล์ สูบบุหรี่ ปื้นขาว ลิ้น"}:

\begin{verbatim}
[+] Test DenseNet-BERT -> HTTP 200
{
  "success": true,
  "model_used": "DenseNet-BERT",
  "result": {
    "disease_name": "Leukoplakia",
    "risk_level": "High",
    "confidence_score": 0.912,
    "bounding_boxes": [],
    "processed_image_url": "/static/outputs/gradcam_densenet_3b91ef.jpg"
  }
}
\end{verbatim}

% ------------------------------------------------------------------------------
% ภาคผนวก ง: พจนานุกรมอาการทางคลินิกและเกณฑ์การประเมินความเสี่ยง
% ------------------------------------------------------------------------------
\chapter{พจนานุกรมอาการทางคลินิกและเกณฑ์การประเมินความเสี่ยง}
\label{appendix:clinical_dictionary}

ภาคผนวกนี้แสดงรายละเอียดชุดคำสำคัญทางคลินิกสำหรับใช้ในการจับคู่กฎเกณฑ์ (Rule Matching) และเกณฑ์การจำแนกระดับความเสี่ยงของรอยโรคที่นำมาใช้ในระบบสนับสนุนการตัดสินใจทางคลินิก

\section{พจนานุกรมคำสำคัญทางคลินิกจำเพาะตามรอยโรค}
ตารางที่ \ref{tab:clinical_symptoms} แสดงรายการคำสำคัญและอาการที่ใช้ตรวจสอบความสอดคล้องระหว่างข้อมูลซักประวัติของผู้ป่วยกับชนิดรอยโรคที่แบบจำลองตรวจพบ:

\begin{table}[htbp]
\centering
\small
\caption{พจนานุกรมคำสำคัญและอาการทางคลินิกจำเพาะตามรอยโรคช่องปาก}
\label{tab:clinical_symptoms}
\begin{tabular}{p{3.2cm}p{5.5cm}p{4.5cm}}
\toprule
\textbf{ชื่อรอยโรค (Pathology)} & \textbf{คำสำคัญภาษาไทยที่ระบบรองรับ} & \textbf{รหัสฟิลด์ข้อมูลในระบบ} \\
\midrule
Leukoplakia & ดื่มเหล้า, สูบบุหรี่, เคี้ยวหมาก, แอลกอฮอล์ & \texttt{drinkalcohol}, \texttt{smoking}, \texttt{chewbetelnut} \\
\midrule
Lichen planus & กินเผ็ด, กินเผ็ดแสบ, รับประทานอาหารเผ็ด & \texttt{eatspicyfood} \\
\midrule
Candidiasis & เช็ดออก, เช็ดออกได้, คราบขาวที่สามารถลอกออก & \texttt{wipeoff} \\
\midrule
Linea alba buccalis & ไม่มีอาการ, ไม่มีอาการผิดปกติ & \texttt{nosymptoms} \\
\midrule
Ulcer & เจ็บเมื่อโดนแผล, มีอาการเจ็บ, เจ็บตลอดเวลา, ระคายเคืองตลอดเวลา & \texttt{alwayshurts} \\
\bottomrule
\end{tabular}
\end{table}

\section{เกณฑ์การจัดระดับความเสี่ยงของรอยโรคทางคลินิก (Risk Stratification)}
ระบบจัดระดับความเสี่ยง (Risk Level) เพื่อช่วยให้ทันตแพทย์สามารถจัดลำดับความเร่งด่วนในการส่งต่อหรือตรวจชิ้นเนื้อเพิ่มเติม (Biopsy) โดยแบ่งออกเป็น 3 ระดับตามตารางที่ \ref{tab:risk_levels}:

\begin{table}[htbp]
\centering
\small
\caption{เกณฑ์การจำแนกระดับความเสี่ยงของรอยโรคทางช่องปาก}
\label{tab:risk_levels}
\begin{tabular}{p{2.5cm}p{4cm}p{3cm}p{4cm}}
\toprule
\textbf{ระดับความเสี่ยง} & \textbf{ชนิดรอยโรคที่เข้าเกณฑ์} & \textbf{เงื่อนไขความมั่นใจ} & \textbf{คำแนะนำในการปฏิบัติทางคลินิก} \\
\midrule
\textbf{ความเสี่ยงสูง\newline(High Risk)} & Leukoplakia, Ulcer, มะเร็งช่องปาก (Oral Cancer) & ค่าคะแนนความเชื่อมั่น $\ge 0.50$ & มีโอกาสพัฒนาเป็นมะเร็งช่องปาก แนะนำให้ตรวจชิ้นเนื้อหรือส่งต่อแพทย์ผู้เชี่ยวชาญ \\
\midrule
\textbf{ความเสี่ยงปานกลาง\newline(Medium Risk)} & Lichen planus, หรือกลุ่มความเสี่ยงสูงที่มีความเชื่อมั่นต่ำกว่าเกณฑ์ & ค่าคะแนนความเชื่อมั่น $\ge 0.60$ & แนะนำให้นัดหมายติดตามอาการอย่างใกล้ชิดภายใน 2-4 สัปดาห์ \\
\midrule
\textbf{ความเสี่ยงต่ำ\newline(Low Risk)} & Candidiasis, Linea alba buccalis, หรือรอยโรคทั่วไป & ทุกระดับความเชื่อมั่น & รอยโรคไม่ใช่มะเร็ง รักษาตามอาการหรือให้คำแนะนำการรักษาสุขอนามัยช่องปาก \\
\bottomrule
\end{tabular}
\end{table}

\section{การกำหนดสีสัญลักษณ์และรูปแบบการแสดงผลทางคลินิก}
เพื่อช่วยให้ทันตแพทย์สามารถสังเกตความผิดปกติได้อย่างรวดเร็ว ระบบได้กำหนดสีของเส้นกรอบ Bounding Box จำเพาะตามประเภทของรอยโรค ดังนี้:
\begin{itemize}
    \item \textbf{Leukoplakia:} เส้นกรอบสีเขียว (Green - \texttt{BGR: 0, 255, 0})
    \item \textbf{Lichen planus:} เส้นกรอบสีแดง (Red - \texttt{BGR: 0, 0, 255})
    \item \textbf{Candidiasis:} เส้นกรอบสีขาว (White - \texttt{BGR: 255, 255, 255})
    \item \textbf{Linea alba buccalis:} เส้นกรอบสีเทา (Gray - \texttt{BGR: 128, 128, 128})
    \item \textbf{Ulcer:} เส้นกรอบสีส้ม (Orange - \texttt{BGR: 255, 140, 0})
\end{itemize}
